\documentclass[10pt,a4paper]{article}
\usepackage[margin=2.2cm]{geometry}
\usepackage{fontspec}
\usepackage{kotex}
\setmainfont{Apple SD Gothic Neo}
\setsansfont{Apple SD Gothic Neo}
\usepackage{amsmath,amssymb}
\usepackage{graphicx}
\usepackage{booktabs}
\usepackage{caption}
\usepackage[dvipsnames]{xcolor}
\usepackage[colorlinks=true,linkcolor=MidnightBlue,citecolor=MidnightBlue,
            urlcolor=MidnightBlue]{hyperref}
\usepackage{titlesec}
\titleformat{\section}{\large\bfseries}{\thesection}{0.6em}{}
\titleformat{\subsection}{\normalsize\bfseries}{\thesubsection}{0.5em}{}
\captionsetup{font=small,labelfont=bf}
\setlength{\parskip}{0.35em}
\setlength{\parindent}{0pt}

% 주장 상자 — 모든 주장은 반증 조건을 함께 적는다
\usepackage{mdframed}
\newmdenv[linewidth=0.4pt,linecolor=gray!60,backgroundcolor=gray!5,
          innertopmargin=6pt,innerbottommargin=6pt,skipabove=8pt,skipbelow=8pt]{claimbox}
\newcommand{\claim}[2]{\begin{claimbox}\textbf{주장.} #1\par\smallskip
  \textcolor{gray!70}{\textbf{반증 조건.} #2}\end{claimbox}}

\title{\bfseries 새 정보인가와 법칙인가는 다른 물음이다\\[2pt]
\large 출시 전 램프와 출시 후 감쇠}
\author{월드모델 실험실 · 노트 655--656}
\date{2026-08-06}
\begin{document}\maketitle

\section{한 줄}

소셜 법칙 트랙에서 특성 둘을 쟀다. 하나는 \textbf{판에 이미 있었고}(재인코딩)
도메인 부호가 6:3 이었다. 다른 하나는 \textbf{판에 없던 새 정보인데} 도메인
부호가 정확히 4:4 다. \textbf{관문을 통과한 것이 법칙임을 뜻하지 않는다.}

\section{자료}

\verb|ingest/wiki_views| 가 노트 149 규약대로 \textbf{오픈 이전 90일}만 긁어
뒀다. 그래서 사용자가 지목한 ``조회수가 올라간다''를 한 번도 못 봤다 --- 성장은
출시 \emph{후}에 있다. 오픈 이후 90일 창을 새로 긁었다(2,330건 중 2,329건이
30일 이상).

\textbf{제목을 다시 안 풀었다.} 기존 캐시에 \verb|page| 가 있고 오픈일은
\emph{이전 창의 마지막 날 $+$ 1} 로 복원된다(창이 $d_0{-}90 \sim d_0{-}1$).

\textbf{이 자료는 축이 될 수 없다.} 오픈 이후 값은 정의상 오픈 시점에 관측
불가다. 그래서 캐시를 다른 디렉토리에 두고, \verb|LAW_ONLY| 상수를 두고,
\textbf{감사가 \texttt{lab/*axes.py} 의 참조를 잡게} 했다 --- 주석으로 적은
금지는 다음 사람이 안 읽는다(노트 598 $\cdot$ 546 $\cdot$ 549).

\section{실행 전에 결함 하나를 잡았다}

캐시 파일에 \verb|lang| 필드가 \textbf{없다}. 그런데 \verb|wiki_views| 는
애니 $\cdot$ 웹툰 $\cdot$ 팝업 $\cdot$ 도서 $\cdot$ 펀딩 $\cdot$ 아이돌을
\textbf{한국어 위키}에서 긁는다. \verb|j.get("lang") or "en"| 으로 뒀으면 그
여섯 도메인이 한국어 제목을 영문 위키에 물어 \textbf{전부 빈 응답}이 된다 ---
``출시 후 자료가 절반뿐이다''로 읽었을 것이다. 레코드 접두로 복원하게 고쳤고,
수집 428건에서 \textbf{빈 응답 0}, 그중 애니 397건이 정상 수집됐다.

\section{두 물음이 갈렸다}

\begin{center}\small
\begin{tabular}{lrr}\toprule
특성 & 기존 축과 최대 $|\rho|$ & 도메인 부호 일치 \\ \midrule
출시 전 램프 & $\mathbf{0.94}$ (재인코딩) & 6:3 \\
출시 후 감쇠율 & $0.62$ \textbf{(새 정보)} & $\mathbf{4:4}$ \\
출시 후 꼬리 두께 & $0.60$ (새 정보) & 4:4 \\
출시 후 반감기 & $0.52$ (새 정보) & 4:4 \\ \bottomrule
\end{tabular}\end{center}

\claim{관문 ③(기존 축과의 상관)을 통과하면 쓸 만한 특성이다.}{통과했는데 도메인
부호가 동전이면 그 특성은 새롭지만 \emph{옮겨 가지 않는다} --- 두 성질은 독립이다.}

출시 전 램프는 \verb|wiki_momentum| 과 $0.89\sim0.94$ 였다. momentum 이 곧
램프이므로 당연하다 --- 그것은 \textbf{판에 이미 있던 것}이었다. 출시 후 감쇠는
최대 $0.62$(게임에서만)로 문턱 $0.85$ 아래다. \textbf{판에 없던 것}이다.

그런데 부호가 도메인마다 뒤집힌다. 감쇠율 $\rho$: 세계애니 $+0.233$ $\cdot$
시장팝업 $+0.220$ $\cdot$ 애니 $+0.137$ $\cdot$ 모바일 $+0.032$ 대
게임 $\mathbf{-0.217}$ $\cdot$ 팝업 $-0.162$ $\cdot$ 만화 $-0.096$ $\cdot$
웹툰 $-0.030$. \textbf{정확히 4:4 다.}

\section{한 자료에서 두 이야기가 나온다}

노트 655 가 조항을 하나 세웠다 --- \emph{일치율은 행 가중과 도메인 수를 같이
적는다}. 이번에 그것이 값을 했다.

\begin{center}\small
\begin{tabular}{lrr}\toprule
모수 & 행 가중 일치 & 도메인 일치 \\ \midrule
감쇠율 & 0.811 & 4/8 \\
꼬리 두께 & 0.811 & 4/8 \\
반감기 & \textbf{0.201} & 4/8 \\ \bottomrule
\end{tabular}\end{center}

세계애니가 2,160행 중 \textbf{1,128행}이다. 행으로 세면 ``0.811 --- 대체로
맞는다''이고 도메인으로 세면 ``동전''이다. 반감기는 더 극단이라 행 가중이
\textbf{0.201} --- 가장 큰 도메인에서 다수 부호가 \emph{틀렸다}. 하나만 적으면
어느 쪽으로든 오해한다.

\section{내 예측이 틀렸다}

사전 등록에 ``출시 후가 더 잘 넘어갈 것''이라 적었다. 근거는 출시 전 램프가
\emph{위키 문서가 언제 서는지}라는 \textbf{우리 관측 절차}의 성질을 섞고 있고
(세계애니 $+0.856$ 대 웹툰 $-0.021$), 출시 후 감쇠는 \emph{사람들이 잊는 속도}
라 세상의 성질에 가까울 것이라는 것이었다.

결과는 $6{:}3 \to 4{:}4$ 로 \textbf{나빠졌다}. 그런데 관문 ③ 은 통과했다.
그러므로 ``절차가 아니라 세상''이라는 진단의 \emph{절반}은 맞았다 --- 새 정보인
것은 맞고, 그것이 옮겨 간다는 것은 틀렸다.

\section{남는 물음}

감쇠율 양수는 \emph{출시 후에도 조회가 는다}는 뜻이다. 애니 $\cdot$ 세계애니
에서는 흥행과 같이 가고 게임 $\cdot$ 팝업에서는 반대다. 게임의 상승하는 위키
곡선이 \textbf{문제 있는 출시가 만든 검색}일 수 있다 --- 그렇다면 같은 모수가
도메인마다 \emph{다른 것을 뜻하는} 것이고, 그건 법칙이 아니라 해석의 문제다.
\textbf{안 쟀다.}

\textbf{처분}: 축으로 안 쓴다(시간 게이트 밖). 법칙 후보로도 채택 안 한다(4/8).
자료는 남긴다 --- 부호가 뒤집히는 이유를 가르면 \textbf{어느 도메인들이 같은
기제인가}의 근거가 된다.

\begin{figure}[h]\centering
\includegraphics[width=\textwidth]{figs/f1.pdf}
\caption{(a) 두 물음의 평면. 오른쪽 아래 녹색이 법칙 후보 구역이고 둘 다
거기 없다 --- 하나는 위(재인코딩), 하나는 왼쪽(전이 없음). (b) 감쇠율의 부호가
4:4 로 갈린다. (c) 같은 자료에서 행 가중과 도메인 수가 반대 이야기를 한다.}
\end{figure}

\end{document}
